%==============================================================
% ICCIT 2026 submission -- IEEE conference format
% Requirements enforced here: A4 paper, max 6 pages, ANONYMOUS.
%
% Build:  pdflatex iccit2026 && pdflatex iccit2026
%
% >>> BEFORE SUBMISSION, resolve every "TODO" comment below. <<<
%==============================================================
\documentclass[conference,a4paper]{IEEEtran}
\IEEEoverridecommandlockouts

\usepackage{cite}
\usepackage{amsmath,amssymb,amsfonts}
\usepackage{graphicx}
\usepackage{booktabs}
\usepackage{textcomp}
\usepackage{url}

\graphicspath{{figures/}}

% TODO(verify): confirm the exact hosted model identifier actually used in
% the experiments. The submitted draft said "gemini-3.5-flash", which is not a
% released model name. Fix this macro once, it propagates everywhere.
\newcommand{\llmProprietary}{gemini-2.5-flash}
\newcommand{\llmOpen}{gpt-oss-120b}

\begin{document}

\title{Multimodal Bangla Fake News Image Detection:\\
Modality Ablation, Distribution-Shift Evaluation,\\
and an Evidence-Grounded Verification Prototype}

%--------------------------------------------------------------
% ANONYMOUS submission block (ICCIT 2026 requires anonymity).
% Do NOT add funding notes or acknowledgments to the blind version.
%--------------------------------------------------------------
\author{\IEEEauthorblockN{Anonymous Submission}
\IEEEauthorblockA{Paper ID: XXXX}}

% --- Camera-ready author block: uncomment only after acceptance ---
% \author{\IEEEauthorblockN{1\textsuperscript{st} Given Name Surname}
% \IEEEauthorblockA{\textit{dept.} \\ \textit{organization} \\ City, Country \\ email}
% \and
% \IEEEauthorblockN{2\textsuperscript{nd} Given Name Surname}
% \IEEEauthorblockA{\textit{dept.} \\ \textit{organization} \\ City, Country \\ email}}

\maketitle

\begin{abstract}
Fake news detection for low-resource languages such as Bangla is constrained
by the scarcity of image-centric training data and by classifiers that report
only same-distribution accuracy. This paper presents a two-stage multimodal
pipeline for manipulated news imagery circulating on Bangladeshi social media.
The first stage fuses a BLIP-generated caption, encoded with multilingual
SBERT and refined by Cross-Image Semantic Fusion (CISF), with a CLIP visual
embedding, and classifies the concatenated vector using gradient-boosted
decision trees. A per-modality ablation on the BanFakeMM benchmark, under
repeated stratified cross-validation, shows that CLIP visual features account
for nearly all of the discriminative signal: CLIP alone reaches 95.58 percent
accuracy, and no text-inclusive configuration exceeds it beyond fold-to-fold
variation. Optical character recognition performs near chance in isolation
(56.80 percent) and is removed from the final pipeline. Evaluated on a
disjoint AdversarialGen manipulation style, zero-shot accuracy falls to
86.0 percent with fake-recall at 73.0 percent, quantifying a substantial
distribution-shift gap; threshold adjustment and supervised adaptation on
target-style data raise accuracy to 91.5 percent, which we report explicitly
as domain adaptation rather than unseen-style generalization. The second stage
is a retrieval-augmented, dual-model verification prototype grounded in a
Bangla fact-check corpus; it is described and characterized qualitatively, and
is not claimed as a validated component.
\end{abstract}

\begin{IEEEkeywords}
Bangla fake news detection, multimodal misinformation, CLIP, XGBoost,
distribution shift, domain adaptation, retrieval-augmented generation,
low-resource languages
\end{IEEEkeywords}

%==============================================================
\section{Introduction}
%==============================================================

Bangladesh's social media ecosystem has become a primary vector for
manipulated news imagery, ranging from digitally altered photocards and forged
documents to AI-generated frames misattributed to public figures. Unlike
text-only misinformation, this class of content exploits the visual and
textual channels jointly: a viral image is rarely judged on its pixels alone,
but on the claim attached to it as it spreads.

The problem is compounded by linguistic specificity. Most fake-news detection
methods have been built and validated on English datasets from Western
platforms, and presuppose signals such as dense comment threads or public
retweet graphs that are sparse or structured differently in the Bangladeshi
context. Three concrete gaps follow. First, no large-scale, organically
collected corpus of manipulated Bangladeshi news images exists, so training
must rely on synthetically constructed benchmarks. Second, an automatically
generated caption describes what is visually present in an image, not the
claim being asserted about it, so a verdict cannot be reached from the image
alone. Third, same-distribution accuracy is routinely reported as if it were
real-world reliability, without a disjoint evaluation that tests whether a
classifier has merely learned one narrow manipulation style.

This paper addresses these gaps with three contributions.
\begin{itemize}
\item A per-modality ablation that determines which modalities actually
contribute, reported with precision, recall and F1 alongside accuracy, and
which leads us to remove OCR from the pipeline rather than include every
available signal by default.
\item A distribution-shift evaluation on a disjoint BanFakeMM manipulation
style that quantifies a 9.4-point zero-shot accuracy gap and a 23-point
fake-recall gap, together with an explicit separation between what is
recovered by threshold adjustment and what is recovered by \emph{supervised
adaptation} on target-style data.
\item A retrieval-augmented dual-model verification stage, grounded in a
Bangla-language fact-check corpus and supporting direct image-to-image
evidence matching, presented and scoped as a \emph{prototype} pending a
labelled quantitative evaluation.
\end{itemize}

We deliberately do not claim architectural novelty for the fusion or
verification paradigms themselves; the regional adaptation and the evaluation
protocol are where we locate this work's contribution.

%==============================================================
\section{Related Work}
%==============================================================

Research on Bangla fake news detection remains limited relative to English.
BanFakeNews \cite{banfakenews} established an early large-scale benchmark of
roughly 50{,}000 annotated Bangla news items. MultiBanFakeDetect
\cite{multibanfake} extended this line to multimodal settings with a hybrid
DenseNet-169/mBERT fusion architecture, and is the closest prior system to the
present work in both language and modality scope. BanFakeNews-2.0
\cite{banfakenews2} contributed a 60{,}000-instance manually curated Bangla
corpus, but is text-only. Beyond Bangla, MMCFND \cite{mmcfnd} extends
caption-aware multimodal detection to seven low-resource Indic languages
including Bengali.

None of these systems is built around a defined image-manipulation task with
adversarially constructed positive examples: MultiBanFakeDetect labels whole
articles by their overall claim, and BanFakeNews-2.0 is text-only. Nor does
any incorporate an external evidence-retrieval stage. Outside the Bangla
literature, EANN \cite{eann}, MVAE \cite{mvae} and SpotFake \cite{spotfake}
established that visual and textual signals are complementary for fake-news
classification, but each was developed on English-language, Western-platform
data and none pairs classification with evidence verification.

A more recent line of work pairs multimodal classification with
retrieval-augmented, multi-model verification directly. DEFAME \cite{defame}
orchestrates specialized LLM-based experts against retrieved multimodal
evidence; RAGAR \cite{ragar} applies retrieval-augmented reasoning to political
fact-checking; RAMA \cite{rama} proposes a retrieval-augmented multi-agent
framework; and CaLM \cite{calm} contrasts a large and a small language model to
verify grounded generation. None is Bangla-specific or grounded in a
Bangladeshi evidence corpus. Correspondingly, we treat the dual-model
verification paradigm as established prior art and claim only its regional
grounding, which is also why we report that component as a prototype rather
than as a validated system.

%==============================================================
\section{Methodology}
%==============================================================

\subsection{System Overview}

The pipeline has two cooperating stages (Fig.~\ref{fig:arch}). The first is a
local multimodal classifier that assigns a fake-probability score to an
uploaded image from its visual and captioned content alone. The second
retrieves independently sourced fact-check evidence---matched either by the
image's pixels or by its associated claim text---and produces an
evidence-based verdict using two independently configured large language
models. The stages are complementary: the classifier provides a fast,
dataset-trained risk signal even when no external evidence exists, while the
verification stage provides claim-level, provenance-backed reasoning whenever
relevant evidence can be found.

\begin{figure}[t]
\centering
\includegraphics[width=\columnwidth]{fig_architecture.png}
\caption{Two-stage pipeline: local multimodal classifier (top) and
retrieval-augmented dual-model verification prototype (bottom).}
\label{fig:arch}
\end{figure}

\subsection{Dataset and Partitioning}

The classifier is trained and evaluated on the AdversarialEdit split of the
public BanFakeMM benchmark \cite{banfakemm}: 2{,}000 images, 1{,}000 authentic
and 1{,}000 manipulated. BanFakeMM constructs its manipulated images through a
defined editing or generation pipeline rather than by subjective annotation of
ambiguous real-world content, so ground-truth labels are deterministic by
construction.

To assess behaviour beyond this distribution, we use BanFakeMM's
AdversarialGen split, a fully AI-generated fake-image collection with no
image-pool overlap with the training data. Of this split, 300 images are
reserved for the supervised-adaptation experiment of Section~\ref{sec:shift}
and a disjoint 200 images (100 authentic, 100 manipulated) form the shifted
evaluation set.

% TODO(reviewer-2, BLOCKING): disjointness is currently established by
% filename comparison only. Before submission, re-verify with a perceptual
% hash (pHash/dHash) AND a CLIP-embedding near-duplicate sweep across all
% partitions, report the number of duplicate clusters removed, and rebuild
% the cross-validation folds with GroupKFold grouped by duplicate cluster and
% source article. Then delete the disclosure sentence below and state the
% verified result instead.
\emph{Limitation.} Partition disjointness is currently established by
filename comparison rather than by perceptual-hash or embedding-based
near-duplicate detection, and cross-validation folds are assigned at the image
level by stratified randomization rather than grouped by source article or
near-duplicate cluster. Residual near-duplicate leakage is therefore an
unresolved risk rather than a ruled-out one, and the same-distribution figures
below should be read with that caveat.

\subsection{Feature Extraction and Fusion}

For each image, three representations are extracted. A BLIP model generates a
natural-language caption, which is encoded into a dense semantic embedding by
a multilingual SBERT model and refined through Cross-Image Semantic Fusion
(CISF): the embedding is combined with the weighted average of its $k=5$
nearest semantic neighbours in the training corpus using fusion weight
$\alpha=0.65$ in favour of the original embedding, then re-normalized and
reduced to 32 principal components. Each training embedding's neighbour search
excludes itself, and at inference an unseen embedding is compared only against
this fixed training-set reference. In parallel, a CLIP vision encoder extracts
a 512-dimensional visual embedding, reduced to 64 principal components by its
own transform. The caption is separately tokenized and stemmed, with separate
handling for Bangla and Latin script, to produce TF-IDF features.

The final vector concatenates four blocks, each fit only on the training fold:
the 32-component CISF caption embedding, the 64-component CLIP embedding, a
TF-IDF representation over the processed caption (unigrams and bigrams,
vocabulary capped at 5{,}000, minimum document frequency 3), and auxiliary
numeric features (token counts, neighbour-similarity statistics).

OCR text extraction was evaluated during ablation and removed from the final
pipeline; Section~\ref{sec:ablation} reports the evidence for that decision.

The values $k=5$, $\alpha=0.65$ and 32 components were fixed design choices
rather than outcomes of a validation search; a sensitivity study over these,
and comparison against adaptive fusion mechanisms such as MMFusionNet
\cite{mmfusionnet} or CLAAF \cite{claaf}, is left for future work.

\subsection{Training Configuration and Reproducibility}

The classifier is an XGBoost gradient-boosted decision tree. Hyperparameters
were selected by randomized search (8 sampled configurations under 3-fold
cross-validation, scored by fake-class F1). The selected configuration is:
\texttt{max\_depth}~3, \texttt{learning\_rate}~0.0742,
\texttt{n\_estimators}~443, \texttt{subsample}~0.866,
\texttt{colsample\_bytree}~0.683, \texttt{min\_child\_weight}~2,
\texttt{gamma}~0.143, \texttt{reg\_alpha}~$1.01\times10^{-5}$,
\texttt{reg\_lambda}~4.911. The search budget is small and was chosen for
computational cost; a wider budget is left for future tuning.

Ablation results come from repeated stratified 5-fold cross-validation with
2 repeats (10 folds total), with identical fold splits and fixed
hyperparameters held constant across every evaluated configuration. All
dimensionality-reduction transforms, vocabularies and neighbour reference sets
are fit inside the training fold only.

% TODO(reviewer-6): fill in the exact values actually used, then delete this
% comment. These are required for the reproducibility paragraph to be true.
\emph{Reproducibility.} All runs use a fixed random seed
(\texttt{seed}~=~[SEED]) for fold assignment, the randomized search and the
classifier. Software versions, exact BLIP/CLIP/SBERT checkpoint identifiers,
hardware and end-to-end runtime are given in the supplementary material, and
code and feature files will be released on acceptance.

\subsection{Retrieval-Augmented Dual-Model Verification (Prototype)}
\label{sec:verify}

The evidence corpus is built independently of the classifier's training data,
from a single Bangla-language fact-checking publisher, via sitemap-based and
keyword-targeted scraping. Of 1{,}521 indexed documents, 1{,}465 (96 percent)
include full body text; the remaining 56 fall back to a headline-level
summary. The single-publisher composition is a scoping limitation, not a claim
of broad coverage.

The claim to be verified is resolved by priority: user-supplied claim text if
provided, otherwise the BLIP caption as a documented fallback. In parallel,
the uploaded image's CLIP embedding is compared against embeddings of images
from previously fact-checked articles. A cosine similarity of 0.90 is treated
as evidence that this exact image has already been independently verified;
during development, unrelated pairs scored 0.55--0.60 while genuine re-shared
copies scored above 0.90. When a direct match is found the system reports the
matched article and its similarity as evidence rather than asserting a truth
verdict, since an image match establishes only that the image has previously
circulated and been checked, not that the claim currently attached to it is
the same claim the matched article addressed.

When no direct match exists, the resolved claim queries the local evidence
store by SBERT semantic similarity (top-5, minimum similarity 0.30), and live
web-search results are retrieved separately. Both are passed to two
independently configured models---\texttt{\llmProprietary} and the open-weight
\texttt{\llmOpen}, at temperature 0---each given an identical prompt and
instructed to return a JSON verdict (\emph{supported}, \emph{contradicted},
\emph{insufficient}) with a confidence score, rationale and cited source URLs.
Each model is queried independently with no visibility into the other's
output. A verdict is returned only when both models complete and agree; any
disagreement or failure is routed to human review rather than resolved by
whichever response is available.

% TODO(reviewer-5, BLOCKING if the prototype framing is dropped): to promote
% this from prototype to evaluated component, build a labelled claim set
% (target 150-200 claims with adjudicated supported/contradicted/insufficient
% labels), then report: retrieval recall@1 and recall@5, verdict accuracy and
% macro-F1, abstention rate, image-match precision at the 0.90 threshold, a
% table of 2-3 worked examples, and a single-model (each model alone) versus
% dual-model consensus comparison. Until those exist, the scoping paragraph
% below must stay.
\emph{Scope.} Both the 0.90 similarity threshold and the consensus rule are
practical operating points set from informal development observation, not from
a calibrated precision/recall study. No labelled claim set, retrieval metric,
verdict-accuracy measurement, or single-versus-dual-model comparison is
reported in this paper, and no sustained deployment trial was conducted.
Accordingly this stage is presented as a working prototype that demonstrates
feasibility and defines an evaluation target, not as a validated verification
system, and none of the quantitative claims in this paper depend on it.

%==============================================================
\section{Experimental Results}
%==============================================================

\subsection{Per-Modality Ablation}
\label{sec:ablation}

\begin{table}[t]
\caption{Per-Modality Ablation, Repeated Stratified 5-Fold CV
($2$ Repeats, $10$ Folds). Accuracy Shown as Mean $\pm$ Standard Deviation
Across Folds. All Values in Percent.}
\label{tab:ablation}
\centering
\setlength{\tabcolsep}{3.5pt}
\begin{tabular}{lcccc}
\toprule
\textbf{Configuration} & \textbf{Acc.} & \textbf{Prec.$_{f}$} &
\textbf{Rec.$_{f}$} & \textbf{F1$_{f}$} \\
\midrule
OCR only              & 56.80 $\pm$ 1.78 & 63.44 & 32.55 & 42.94 \\
Caption only          & 74.40 $\pm$ 1.79 & 75.05 & 73.15 & 74.07 \\
Text only             & 76.55 $\pm$ 2.01 & 76.75 & 76.25 & 76.48 \\
Text + CISF           & 77.48 $\pm$ 1.37 & 77.43 & 77.65 & 77.50 \\
\midrule
CLIP only             & \textbf{95.58} $\pm$ 0.92 & 95.12 & 96.10 & 95.59 \\
Text + CLIP           & 95.43 $\pm$ 0.83 & 95.16 & 95.75 & 95.44 \\
CLIP + CISF           & 95.58 $\pm$ 1.01 & 95.25 & 95.95 & 95.59 \\
Caption + CISF + CLIP & 95.40 $\pm$ 0.97 & 95.10 & 95.75 & 95.41 \\
\bottomrule
\multicolumn{5}{l}{\footnotesize Subscript $f$ denotes the manipulated
(``fake'') class.}
\end{tabular}
\end{table}

Table~\ref{tab:ablation} reports every feature configuration under an
identical protocol. Three results matter.

First, OCR text in isolation reaches 56.80 percent accuracy, close to the
50 percent chance baseline for this balanced binary task, with fake-recall of
only 32.55 percent. This is the empirical basis for removing OCR from the
final pipeline. We do not generalize this beyond BanFakeMM: MultiBanFakeDetect
\cite{multibanfake} indicates that text-in-image remains a meaningful vector
for Bangla misinformation more broadly, and BanFakeMM's manipulations are
photographic rather than text-embedded.

Second, all text-only configurations saturate below 78 percent, while any
configuration containing CLIP exceeds 95 percent. Visual features, not caption
text or its derived properties, carry the discriminative signal for this task.

Third---and this is a negative result we report rather than obscure---no
text-inclusive configuration exceeds CLIP alone. CLIP-only reaches
95.58 $\pm$ 0.92 percent; adding CISF on top leaves it unchanged at
95.58 $\pm$ 1.01; the full caption+CISF+CLIP configuration reaches
95.40 $\pm$ 0.97, marginally \emph{below} CLIP alone. The largest gap among
CLIP-containing configurations is 0.18 points against a per-fold standard
deviation of roughly 0.9--1.0 points, so these configurations are
indistinguishable at this sample size.

% TODO(reviewer-4, BLOCKING): the fusion architecture is not currently
% justified by these numbers. Before submission do ONE of:
%   (a) demonstrate a measurable benefit -- run CLIP-only vs. fusion on the
%       shifted set of Table II, report a paired test (5x2cv F-test or
%       bootstrap 95% CI on the difference), and add calibration (ECE/Brier)
%       and AUROC; if fusion wins on robustness or calibration, say so and
%       keep it; OR
%   (b) simplify -- make CLIP + XGBoost the proposed model, retitle, and
%       demote CISF to a reported negative result.
% The paragraph below is honest but is NOT a substitute for either.
\emph{On the fusion design.} On same-distribution accuracy the caption
pathway earns its place only if it contributes robustness or calibration
rather than accuracy, which Table~\ref{tab:ablation} does not measure.
Retaining a multimodal architecture is not by itself a justification. We
therefore report the final configuration's parity with CLIP-only as an open
question about the fusion design rather than as a positive result, and treat
CLIP-only as the operative baseline any fusion must beat.

\subsection{Evaluation Under Distribution Shift}
\label{sec:shift}

\begin{table}[t]
\caption{Evaluation on the Disjoint AdversarialGen Manipulation Style
(200 Images, 100/100 Balanced). All Values in Percent.}
\label{tab:shift}
\centering
\setlength{\tabcolsep}{3pt}
\begin{tabular}{llccc}
\toprule
\textbf{Setting} & \textbf{$\tau$} & \textbf{Acc.} &
\textbf{Rec.$_{f}$} & \textbf{Rec.$_{r}$} \\
\midrule
\multicolumn{5}{l}{\emph{Zero-shot transfer (no target-style data in training)}} \\
Default threshold              & 0.50 & 86.0 & 73.0 & 99.0 \\
Adjusted threshold$^{\dagger}$ & 0.30 & 90.0 & 86.0 & 94.0 \\
\midrule
\multicolumn{5}{l}{\emph{Supervised domain adaptation ($+300$ AdversarialGen images)}} \\
Adapted model$^{\dagger}$      & 0.30 & 91.5 & 91.0 & 92.0 \\
\bottomrule
\multicolumn{5}{p{0.92\columnwidth}}{\footnotesize $^{\dagger}$Threshold
selected by sweeping this same evaluation set. These two rows are optimistic
and are not held-out estimates; see the disclosure in the text. Subscripts $f$
and $r$ denote the manipulated and authentic classes.}
\end{tabular}
\end{table}

Table~\ref{tab:shift} reports the shifted evaluation. The headline result is
the first row: at the default threshold, accuracy falls from 95.40 percent to
86.0 percent, and fake-recall falls from 95.75 percent to 73.0 percent. The
model misses 27 of 100 manipulated images while misclassifying only 1 of 100
authentic ones---a strongly asymmetric error pattern consistent with a
decision boundary calibrated tightly around the training distribution's
specific manipulation signature. The 23-point fake-recall collapse, not the
9.4-point accuracy drop, is the operationally meaningful number: fake-recall
is the safety-critical metric here, since a missed manipulated image
circulates unchecked whereas a false alarm costs only a review step.

Two mitigations are reported, and we separate them deliberately. Lowering the
threshold to 0.30 without retraining raises accuracy to 90.0 percent and
fake-recall to 86.0 percent. Retraining with 300 additional AdversarialGen
images, evaluated at the same threshold with the 200-image set held out of
training, raises accuracy to 91.5 percent with a far more balanced matrix
(91 of 100 manipulated images caught, 92 of 100 authentic images passed),
shown in Fig.~\ref{fig:cm}.

\emph{This second result is domain adaptation, not generalization to an unseen
style.} Once target-style images enter training, the evaluation no longer
measures zero-shot transfer, and the 91.5 percent figure must not be read as
evidence that the generalization gap has closed. Only the 86.0 percent row
measures unseen-style behaviour. The corresponding claim in an earlier draft
of this work---that diversification ``narrows the out-of-distribution gap to
91.5 percent''---conflates the two regimes and is withdrawn.

% TODO(reviewer-1, BLOCKING): remove the threshold-selection leakage.
% Re-partition AdversarialGen at the near-duplicate-cluster level into three
% disjoint subsets: ~300 adaptation-training / ~150 validation / ~200 sealed
% test. Fix the selection criterion in advance (argmax fake-class F1 on the
% validation partition), select tau there, then evaluate the sealed test
% partition EXACTLY ONCE. Replace the two daggered rows and delete the
% disclosure paragraph below. Expect the numbers to fall; that is the point.
\emph{Disclosure: threshold selection.} The threshold of 0.30 was selected by
sweeping the same 200-image set on which the daggered rows are reported. Those
two rows are therefore an optimistic upper bound on what threshold adjustment
achieves, not held-out estimates, and the true held-out values are expected to
be lower. Only the zero-shot default-threshold row (86.0 percent) is free of
this selection effect, and it is the only row we rely on for the paper's
claims about distribution shift. A protocol using a validation partition
disjoint from a sealed test partition is required to report the remaining rows
as measurements, and is the immediate next step for this work.

\begin{figure}[t]
\centering
\includegraphics[width=\columnwidth]{fig3_confusion_matrices.png}
\caption{Confusion matrices on the shifted evaluation set: zero-shot at the
default threshold (left) versus the domain-adapted model (right). The
zero-shot error pattern is strongly asymmetric, missing 27 of 100 manipulated
images while misclassifying 1 of 100 authentic ones.}
\label{fig:cm}
\end{figure}

\subsection{Baseline Comparison}

\begin{table}[t]
\caption{Text-Only Baselines Under the Identical Cross-Validation Protocol.
All Values in Percent.}
\label{tab:baselines}
\centering
\setlength{\tabcolsep}{3.5pt}
\begin{tabular}{lcccc}
\toprule
\textbf{Model} & \textbf{Acc.} & \textbf{Prec.$_{f}$} &
\textbf{Rec.$_{f}$} & \textbf{F1$_{f}$} \\
\midrule
Linear SVM (TF-IDF)     & 76.30 $\pm$ 1.96 & 75.31 & 78.30 & 76.77 \\
XGBoost (TF-IDF + num.) & 78.95 $\pm$ 1.35 & 79.06 & 78.80 & 78.91 \\
\midrule
CLIP only               & 95.58 $\pm$ 0.92 & 95.12 & 96.10 & 95.59 \\
Proposed (final)        & 95.40 $\pm$ 0.97 & 95.10 & 95.75 & 95.41 \\
\bottomrule
\end{tabular}
\end{table}

Table~\ref{tab:baselines} compares against two text-only baselines under the
same protocol, with all four metrics reported for every row. Both baselines
sit 16--19 points below any CLIP-containing configuration, confirming that
visual features provide the primary discriminative signal. We include the
CLIP-only row in this table deliberately: it is the baseline the proposed
configuration must beat, and at present it does not.

%==============================================================
\section{Discussion and Limitations}
%==============================================================

\subsection{Why Visual Features Dominate}

The ablation outcome follows from the nature of BanFakeMM's manipulations
rather than from any inherent weakness of text features. The dataset's
manipulated images are produced by pixel-level editing or full-image
generation, so evidence of manipulation is encoded in visual texture, lighting
consistency and generation artifacts---exactly the signal CLIP's encoder is
trained to capture. A caption- or OCR-derived feature can reflect manipulation
only indirectly, which explains both OCR's near-chance standalone performance
and CLIP's dominance once included. This also bounds the finding: on a corpus
where manipulation is carried by embedded text, the ranking could invert.

CISF's lack of measurable contribution on top of CLIP is consistent with the
same explanation. CISF sharpens an embedding by borrowing signal from similar
training items, a mechanism suited to text embeddings that benefit from
smoothing. CLIP's visual embedding already encodes a rich description directly
from pixels, leaving comparatively little redundant noise for neighbour fusion
to average away. We note this is inferred from the ablation pattern rather
than independently verified.

\subsection{Limitations}

Several limitations bound this work, beyond the two disclosures already stated
in Sections~\ref{sec:verify} and \ref{sec:shift}.

Ground-truth labels derive from BanFakeMM's deterministic construction
process, so the classifier's demonstrated competence is specific to detecting
this class of algorithmic image manipulation---not to adjudicating whether a
news item's overall claim is true. The shifted evaluation, while disjoint from
training data, is drawn from a second split of the same underlying benchmark;
generalization to manipulation styles and image sources entirely outside
BanFakeMM, including organically collected Bangladeshi social media imagery,
remains open.

Threshold-independent metrics---AUROC, FPR@95, and calibration error---were
not computed, and a single confusion matrix for the same-distribution split is
unavailable because per-fold counts were not retained under the
cross-validation protocol. Near-duplicate leakage across folds is not ruled
out. Architecture-level alternatives to the data-centric mitigation used
here---CLIPood's fine-tuning approach \cite{clipood} or domain-specific prompt
tuning \cite{dpod}---were not evaluated. Finally, this pipeline addresses
static imagery; extending it to video-based misinformation would require a
substantially different approach.

%==============================================================
\section{Conclusion}
%==============================================================

Practical Bangla misinformation detection demands more than a
same-distribution accuracy figure. This paper reports a multimodal classifier
for manipulated Bangladeshi news imagery together with an evaluation protocol
designed to expose, rather than conceal, where that classifier fails.

Three findings stand. Modality ablation shows CLIP visual features carry
essentially all of the discriminative signal on BanFakeMM, which justified
removing OCR and which leaves the caption-fusion pathway currently unjustified
on accuracy grounds---a negative result we report rather than obscure.
Evaluation on a disjoint manipulation style shows a substantial zero-shot gap,
with fake-recall falling from 95.75 to 73.0 percent, demonstrating that
same-distribution accuracy is a poor proxy for real-world reliability.
Supervised adaptation on target-style data recovers much of that gap, but by
construction measures domain adaptation rather than unseen-style
generalization, and we label it accordingly.

The immediate next steps follow directly from the disclosures above: a sealed
test partition with threshold selection confined to a disjoint validation
split; near-duplicate verification by perceptual hashing and embedding
similarity with group-aware folds; a paired significance and calibration
comparison to decide whether the fusion pathway is retained or the classifier
simplified to CLIP-only; and a labelled claim set with retrieval and verdict
metrics to promote the verification stage from prototype to evaluated
component.

%==============================================================
\begin{thebibliography}{00}

% TODO(reviewer-6, BLOCKING): the BanFakeMM record below could not be
% independently verified -- supply curator, year and DOI from the actual
% Mendeley Data landing page before submission.
\bibitem{banfakemm} ``BanFakeMM: Bangla multimodal fake news dataset,''
Mendeley Data. [Online]. Available: \url{https://data.mendeley.com/}

\bibitem{banfakenews} M. Z. Hossain, M. A. Rahman, M. S. Islam, and S. Kar,
``BanFakeNews: A dataset for detecting fake news in Bangla,'' in
\emph{Proc. 12th Lang. Resources and Evaluation Conf. (LREC)}, Marseille,
France, 2020, pp. 2862--2871.

\bibitem{multibanfake} F. T. J. Faria \emph{et al.},
``MultiBanFakeDetect: Integrating advanced fusion techniques for multimodal
detection of Bangla fake news in under-resourced contexts,''
\emph{Int. J. Inf. Manage. Data Insights}, 2025.

\bibitem{banfakenews2} H. M. Shibu, S. Datta, M. S. Miah, N. Sami,
M. S. Chowdhury, and M. S. Islam, ``From scarcity to capability: Empowering
fake news detection in low-resource languages with LLMs,'' in
\emph{Proc. Conf. Indo NLP (IndoNLP)}, 2025, arXiv:2501.09604.

\bibitem{mmcfnd} S. Bansal, N. S. Singh, S. S. Dar, and N. Kumar,
``MMCFND: Multimodal multilingual caption-aware fake news detection for
low-resource Indic languages,'' 2024, arXiv:2410.10407.

\bibitem{eann} Y. Wang \emph{et al.}, ``EANN: Event adversarial neural
networks for multi-modal fake news detection,'' in \emph{Proc. 24th ACM
SIGKDD Int. Conf. Knowl. Discovery Data Mining}, 2018, pp. 849--857.

\bibitem{mvae} D. Khattar, J. S. Goud, M. Gupta, and V. Varma,
``MVAE: Multimodal variational autoencoder for fake news detection,'' in
\emph{Proc. World Wide Web Conf. (WWW)}, San Francisco, CA, 2019,
pp. 2915--2921.

\bibitem{spotfake} S. Singhal, R. R. Shah, T. Chakraborty, P. Kumaraguru, and
S. Satoh, ``SpotFake: A multi-modal framework for fake news detection,'' in
\emph{Proc. IEEE 5th Int. Conf. Multimedia Big Data (BigMM)}, Singapore, 2019,
pp. 39--47.

\bibitem{defame} T. Braun, M. Rothermel, M. Rohrbach, and A. Rohrbach,
``DEFAME: Dynamic evidence-based fact-checking with multimodal experts,''
2024, arXiv:2412.10510.

\bibitem{ragar} M. A. Khaliq, P. Chang, M. Ma, B. Pflugfelder, and
F. Mileti\'{c}, ``RAGAR, your falsehood radar: RAG-augmented reasoning for
political fact-checking using multimodal large language models,'' in
\emph{Proc. 7th FEVER Workshop, ACL}, 2024, arXiv:2404.12065.

% TODO(reviewer-6, BLOCKING): confirm the RAMA author list from the arXiv
% record before submission; it could not be independently verified.
\bibitem{rama} ``RAMA: Retrieval-augmented multi-agent framework for
misinformation detection in multimodal fact-checking,'' 2025,
arXiv:2507.09174.

\bibitem{calm} I.-H. Hsu \emph{et al.}, ``CaLM: Contrasting large and small
language models to verify grounded generation,'' in \emph{Findings Assoc.
Comput. Linguistics: ACL 2024}, arXiv:2406.05365.

\bibitem{mmfusionnet} J. He, T. Liu, J. Zhao, and B. Turner,
``MMFusionNet: Context-aware dynamic fusion for multi-modal fake news
detection with large vision-language models,'' 2025, arXiv:2508.05684.

\bibitem{claaf} G. Mu, C. Chen, X. Li, Y. Chen, J. Dai, and J. Li,
``CLAAF: Multimodal fake information detection based on contrastive learning
and adaptive agg-modality fusion,'' \emph{PLOS ONE}, May 2025.

\bibitem{clipood} Y. Shu, X. Guo, J. Wu, X. Wang, J. Wang, and M. Long,
``CLIPood: Generalizing CLIP to out-of-distributions,'' in \emph{Proc. 40th
Int. Conf. Machine Learning (ICML)}, 2023, arXiv:2302.00864.

\bibitem{dpod} A. Bhattacharya, D. Brahma, S. Nagaje, A. Asati, V. Verma, and
S. Biswas, ``Can out-of-domain data help to learn domain-specific prompts for
multimodal misinformation detection?'' in \emph{Proc. IEEE/CVF Winter Conf.
Appl. Comput. Vis. (WACV)}, 2025, arXiv:2311.16496.

\end{thebibliography}

\end{document}
